
Neural networks trained with empirical risk minimization learn spurious correlations, failing systematically on minority groups where spurious features conflict with labels. This work investigates whether training dynamics contain detectable signatures that distinguish minority from majority samples without requiring group annotations. On the Waterbirds dataset, minority-group samples exhibit 8.$3\times$ higher mean loss at epoch 5 (5\% of training) compared to majority samples (0.158 vs. 0.019, Mann-Whitney U = 705,391, p = 7.36 $\times 10^{-15})$. Using linear probes on frozen ResNet-18 representations, we confirm the mechanism: representations encode spurious features (background) with 91.2\% accuracy versus 80.5\% for core features (bird type) at epoch 5, a 10.7 percentage point gap. Loss-based detection at the 95th percentile threshold achieves 32.9\% precision---6.$6\times$ improvement over the 5\% random baseline---though this precision is fundamentally bounded by the minority base rate, not signal weakness. Contrary to initial hypotheses, with ImageNet-pretrained features both spurious and core feature probe accuracies peak at the same epoch (81), indicating that simplicity bias manifests as an accuracy magnitude gap rather than a temporal gap under transfer learning conditions.

